% исправления 27.05.2005
\section{Игра Эренфойхта}
        \label{ehrenfeucht-game}%
Вернёмся от алгебры к логике и сформулируем общий критерий
элементарной эквивалентности двух интерпретаций некоторой
сигнатуры. Будем считать, что наша сигнатура содержит только
предикатные символы. (Это ограничение не очень существенно, так
как функцию $f(x_1,\dots,x_n)$ можно заменить предикатом
$f(x_1,\dots,x_n)\hm=y$, имеющим на один аргумент больше.) Кроме
того, будем считать, что сигнатура конечна (в некоторый момент
наших рассуждений это будет существенно).

Критерий будет сформулирован в терминах некоторой игры,
называемой \emph{игрой Эренфойхта}.%
\index{Игра Эренфойхта|defin}%
\glossary{\Эренфойхт}
В ней участвуют два игрока,
называемые Новатором (\Н) и Консерватором (\К). Игра определяется
        \index{Новатор (в игре Эренфойхта)}%
        \index{Консерватор (в игре Эренфойхта)}%
выбранной парой интерпретаций; как мы докажем, интерпретации
элементарно эквивалентны тогда и только тогда, когда \К\ имеет
выигрышную стратегию в этой игре.

В начале игры Новатор объявляет натуральное число~$k$. Далее они
ходят по очереди, начиная с \Н; каждый из игроков делает $k$
ходов, после чего определяется победитель.

На $i$\д м ходу \Н\ выбирает элемент в одной из интерпретаций (в
любой из двух, причём выбор может зависеть от номера хода) и
помечает его числом $i$. В ответ \К\ выбирает некоторый элемент
из другой интерпретации и также помечает его числом $i$. После
$k$ ходов игра заканчивается. При этом в каждой интерпретации
$k$ элементов оказываются помеченными числами от $1$ до $k$ (мы
не учитываем, кто именно из игроков их пометил). Обозначим эти
элементы $a_1,a_2,\ldots,a_k$ (для первой интерпретации;
элемент~$a_i$ имеет пометку~$i$) и $b_1,b_2,\ldots,b_k$ (для
второй). Элементы $a_i$ и $b_i$ (с одним и тем же~$i$) будем
называть \emph{соответствующими} друг другу. Посмотрим, найдётся
ли предикат сигнатуры, который различает помеченные элементы
первой и второй интерпретации (то есть истинен на некотором
наборе помеченных элементов в одной интерпретации, но ложен на
соответствующих элементах другой). Если такой предикат найдётся,
то выигрывает Новатор, в противном случае\т Консерватор.

Прежде чем доказывать, что эта игра даёт критерий
элементарной эквивалентности, рассмотрим несколько простых
примеров.

\begin{itemize}
        \item
Среди элементов $a_1,\dots,a_k$ и $b_1,\dots,b_k$ могут быть
одинаковые. Если в нашей сигнатуре есть предикат равенства и в
обеих интерпретациях он интерпретируется как совпадение
элементов, то Консерватор обязан повторять ходы, если их
повторил Новатор (скажем, если $a_i=a_j$, а $b_i\ne b_j$, то
Новатор выигрывает, поскольку предикат равенства истинен в одной
интерпретации, но ложен на соответствующих элементах другой).
        \item
Если интерпретации изоморфны, то у Консерватора есть очевидный
способ выиграть: изоморфизм заранее группирует все элементы в
пары соответствующих. (Это согласуется с тем, что изоморфные
интерпретации элементарно эквивалентны.)
        \item
Рассмотрим сигнатуру упорядоченных множеств (предикаты $=$ и
$<$) и её естественные интерпретации в $\bbN$ и $\bbZ$. Они не
являются элементарно эквивалентными, поскольку среди натуральных
чисел есть наименьшее, а среди целых\т нет. Покажем, что в игре
Эренфойхта для этих интерпретаций выигрывает Новатор.

\Н\ объявляет, что игра будет проведена в два хода и на первом
ходу помечает число $0$ из интерпретации~$\bbN$. В ответ
\К\ вынужден
пометить некоторое число $m$ из $\bbZ$. На втором ходу \Н\
помечает в $\bbZ$ некоторое число, меньшее $m$ (например,
$m-1$). Теперь \К\ проигрывает при любом ответном ходе,
поскольку пометить число, меньшее нуля, он не может.
        \item
Для той же сигнатуры рассмотрим интерпретации в $\bbZ$ и $\bbQ$.
Эти интерпретации не элементарно эквивалентны, поскольку порядок
на рациональных числах плотен, а на целых\т нет. Покажем, что в
игре Эренфойхта снова выигрывает Новатор.

Игра будет проходить в три хода. На первых двух ходах \Н\
помечает числа $0$ и $1$ из $\bbZ$. \К\ должен пометить
некоторые элементы $b_1$ и $b_2$ из $\bbQ$. При этом должно быть
$b_1<b_2$ (иначе \Н\ заведомо выиграет). Тогда на третьем ходу \Н\
помечает любое рациональное число, лежащее строго между $b_1$ и
$b_2$. Так как между $0$ и $1$ нет натуральных чисел, \К\ не
может соблюсти требования игры и проигрывает при любом ходе.
        \item
\label{z-equivalent-z-plus-z-2}%
Рассмотрим теперь упорядоченные множества $\bbZ$ и $\bbZ+\bbZ$.
        \index{Z+Z@$\bbZ+\bbZ$}%
        \index{Элементарная эквивалентность $\bbZ$ и $\bbZ+\bbZ$}%
Как мы видели (с.~\pageref{z-equivalent-z-plus-z}),
они элементарно эквивалентны, и потому должна существовать
выигрышная стратегия для Консерватора. Как же он должен играть?
Кажется разумным поддерживать одинаковые расстояния между
соответствующими элементами в $\bbZ$ и $\bbZ+\bbZ$. Проблема в
том, что в $\bbZ+\bbZ$ некоторые
расстояния бесконечны (между элементами
разных слагаемых). Что же делать Консерватору, если
Новатор пометил два таких элемента?

К счастью для \К, он знает заранее, сколько ходов осталось до
конца игры. Ясно, что если игра скоро кончится, то \Н\ не
удастся отличить бесконечное расстояние от достаточно большого.
Более точно, если до конца игры остаётся $s$ ходов, то \К\ может
считать все расстояния, большие или равные~$2^s$, бесконечно
большими. В конце (при $s=0$) это означает, что все ненулевые
расстояния отождествляются (что правильно, так как в конце важен
лишь порядок). Таким образом, \К\ стремится поддерживать такой
\лк инвариант\пк\ (как сказали бы программисты): соответствующие
элементы в~$A$ и в~$B$ идут в одном и том же порядке, и
расстояния между соответствующими парами соседей одинаковы (при
этом все бесконечно большие расстояния считаются одинаковыми).
Ясно, что такая стратегия гарантирует ему выигрыш; надо лишь
проверить, что поддержать инвариант можно.

При очередном ходе \Н\ возможны несколько случаев. \Н\ мог
разбить \лк конечный\пк\ (меньший $2^s$, где $s$\т число
оставшихся ходов) промежуток на две части. В этом случае
соответствующий промежуток в другом множестве также \лк
конечен\пк\ и имеет ту же длину, так что \К\ должен лишь выбрать
элемент на том же расстоянии от концов. Пусть \Н\ разбил \лк
бесконечный\пк\ (длины $2^s$ или больше) промежуток на две
части. Тогда хотя бы одна из частей будет иметь длину $2^{s-1}$
или больше, то есть на следующем шаге будет считаться \лк
бесконечной\пк. Если обе части \лк бесконечны\пк\ (с точки
зрения следующего шага), то \К\ должен разбить \лк
бесконечный\пк\ (длины $2^s$ или более) промежуток другого
множества на две \лк бесконечные\пк\ (длины $2^{s-1}$ или более)
части; это, очевидно, возможно. Если одна часть \лк
бесконечна\пк, а другая \лк конечна\пк, то надо отложить то же
\лк конечное\пк\ расстояние в другом множестве. Наконец, обе
части не могут быть \лк конечными\пк\ (если каждая меньше
$2^{s-1}$, то в сумме будет меньше $2^s$).

Наконец, новый элемент мог быть больше (или меньше) всех уже
отмеченных элементов интерпретации; в этом случае \К\ должен
отметить элемент другой интерпретации, находящийся на том же
расстоянии от наибольшего (наименьшего) отмеченного элемента
(или на \лк бесконечном\пк\ расстоянии, если такова была
ситуация с выбранным \Н\ элементом).
        %
\end{itemize}
\medskip

\begin{problem}
        %
Кто выигрывает в игре Эренфойхта для
упорядоченных множеств
        (\textbf{а})~%
$\bbZ$ и $\bbR$;
        (\textbf{б})~%
$\bbR$ и $\bbQ$;
        (\textbf{в})~%
$\bbN$ и $\bbN+\bbZ$?
        %
Как он должен играть?
        %
\end{problem}

Приведённые примеры делают правдоподобной связь между наличием
формулы, различающей интерпретации, и выигрышной стратегии для
\Н. При этом число ходов, которое понадобится Новатору,
соответствует \emph{кванторной глубине}%
\index{Кванторная глубина формулы|defin}%
\index{Формула!кванторная глубина|defin}
различающей
интерпретации формулы. Кванторная глубина формулы определяется
так:

\begin{itemize}
        \item
Глубина атомарных формул равна нулю.
        \item
Глубина формул $\varphi\vee\psi$ и $\varphi\wedge\psi$ равна
максимуму глубин формул $\varphi$ и $\psi$.
        \item
Глубина формулы $\neg\varphi$ равна глубине формулы $\varphi$.
        \item
Глубина формул $\exists\xi\,\varphi$ и $\forall\xi\,\varphi$ на
единицу больше глубины формулы $\varphi$.
        %
\end{itemize}

Другими словами, глубина формулы\т это наибольшая \лк глубина
вложенности\пк\ кванторов (максимальная длина цепочки вложенных
кванторов).

Рассмотрим позицию, которая складывается в игре после $k$
ходов \Н\ и \К\ (перед очередным ходом \Н) и
за~$l$~ходов до конца игры (таким образом,
общая длина игры есть $k+l$). В этот момент в каждой из
интерпретаций совместными усилиями \Н\ и \К\ выбрано по $k$
элементов. Пусть это будут элементы $a_1,\dots,a_k$ в одной
интерпретации (назовём её~$A$) и $b_1,\dots,b_k$ в другой ($B$).

\textsf{Лемма}. Если есть формула глубины
$l$ с параметрами $x_1,\dots,x_k$, отличающая $a_1,\dots,a_k$
от $b_1,\dots,b_k$, то в указанной позиции \Н\ имеет выигрышную
стратегию; в противном случае её имеет \К.

Поясним смысл условия леммы. Пусть $\varphi$\т формула глубины
$l$, все параметры которой содержатся в списке $x_1,\dots,x_k$.
Тогда имеет смысл ставить вопрос о её истинности в
интерпретации $A$ при значениях параметров $a_1,\dots,a_k$, а
также в интерпретации $B$ при значениях параметров
$b_1,\dots,b_k$. Если окажется, что в одном случае формула $\varphi$
истинна, а в другом ложна, то мы говорим, что $\varphi$ отличает
$a_1,\dots,a_k$ от $b_1,\dots,b_k$.

Пусть такая формула $\varphi$ существует. Она представляет собой
логическую (бескванторную) комбинацию некоторых формул вида
$\forall \xi \,\psi$ и $\exists\xi \psi$, где $\psi$\т формула
глубины $l-1$. Хотя бы одна из формул, входящих в эту
комбинацию, должна также
отличать
$a_1,\dots,a_k$ от $b_1,\dots,b_k$.
Переходя к отрицанию, можно считать, что эта формула начинается
с квантора существования. Пусть формула $\varphi$, имеющая вид
        $$
\exists x_{k+1} \psi(x_1,\dots,x_k,x_{k+1}),
        $$
истинна для $a_1,\dots,a_k$ и ложна для $b_1,\dots,b_k$. Тогда
найдётся такое $a_{k+1}$, для которого в $A$ истинно
        $$
\psi(a_1,\dots,a_k,a_{k+1}).
        $$
Это $a_{k+1}$ и будет выигрывающим ходом Новатора; при любом
ответном ходе $b_{k+1}$ Консерватора формула
        $$
\psi(b_1,\dots,b_k,b_{k+1})
        $$
будет ложной. Таким образом, некоторая формула глубины $l-1$
отличает $a_1,\dots,a_k,a_{k+1}$ от $b_1,\dots,b_k,b_{k+1}$; 
рассуждая по индукции, мы можем считать, что в
оставшейся $(l-1)$\д ходовой игре \Н\ имеет выигрышную
стратегию. (В конце концов мы придём к ситуации, когда некоторая
бескванторная формула отличает $k+l$ элементов в~$A$ от
соответствующих элементов в~$B$, то есть \Н\ выиграет.)

Обратное рассуждение (если наборы не отличимы никакой формулой
глубины $l$, то \К\ имеет выигрышную стратегию в оставшейся
$l$\д ходовой игре) чуть более сложно. Здесь
важно, что по существу есть лишь конечное число различных формул
глубины $k$.

Точнее говоря, будем называть две формулы (с параметрами)
эквивалентными,
если они одновременно истинны или ложны в любой интерпретации на
любой оценке.
Поскольку сигнатура конечна, существует лишь
конечное число атомарных формул, все параметры которых
содержатся среди $u_1,\dots,u_s$. Существует лишь конечное число
булевых функций с данным набором аргументов, поэтому существует
лишь конечное число неэквивалентных бескванторных формул, все
параметры которых содержатся среди $u_1,\dots,u_s$. Отсюда
следует, что существует лишь конечное число неэквивалентных
формул вида
        $$
\exists u_s\, \psi(u_1,\dots,u_{s}),
        $$
и потому лишь конечное число неэквивалентных формул глубины~$1$,
параметры которых содержатся среди $u_1,\dots,u_{s-1}$. (Здесь
мы снова используем утверждение о конечности числа булевых
функций с данным конечным списком
аргументов, а также возможность
переименовывать переменную под квантором, благодаря которой мы
можем считать, что эта переменная есть $u_s$.) Продолжая эти
рассуждения, мы заключаем, что для любого $l$ и для любого
набора переменных $u_1,\dots,u_n$ существует лишь конечное число
неэквивалентных формул глубины $l$, все параметры которых
содержатся среди $u_1,\dots,u_n$. (Здесь мы существенно
используем конечность сигнатуры.)

Вернёмся к игре Эренфойхта. Пусть
элементы $a_1,\dots,a_k$ нельзя отличить от
элементов $b_1,\dots,b_k$ с помощью
формул глубины~$l$. Опишем выигрышную стратегию для
\К. Пусть \Н\ выбрал произвольный элемент в одной из
интерпретаций, скажем,~$a_{k+1}$. Рассмотрим все формулы глубины
$l-1$ с~$k+1$~параметрами (с точностью до эквивалентности их
конечное число); некоторые из них будут истинны на
$a_1,\dots,a_{k+1}$, а некоторые ложны. Тогда формула, утверждающая
существование $a_{k+1}$ с ровно такими свойствами (после квантора
существования идёт конъюнкция всех истинных формул и отрицаний
всех ложных) будет формулой глубины~$l$, истинной на
$a_1,\dots,a_k$. По предположению эта формула должна быть
истинной и на $b_1,\dots,b_k$, и потому существует $b_{k+1}$ с
теми же свойствами, что и $a_{k+1}$. Этот элемент $b_{k+1}$ и
должен пометить \К. Теперь предположение индукции позволяет
заключить, что в возникшей позиции (где до конца игры $l-1$
ходов) у \К\ есть выигрышная стратегия.

Лемма доказана. Её частным случаем является обещанный
критерий элементарной эквивалентности:

\begin{theorem}
        \label{eurenfeucht}%
Интерпретации $A$ и $B$ элементарно эквивалентны тогда и только
тогда, когда в соответствующей игре Эренфойхта выигрывает
Консерватор.
        %
\end{theorem}

\begin{problem}
        %
Покажите, что условие конечности сигнатуры существенно (без него
из элементарной эквивалентности не следует существование выигрышной
стратегии для \К).
        %
\end{problem}

Заметим, что в некоторых случаях (например, для
$\bbZ$ и $\bbZ+\bbZ$) игра Эренфойхта даёт нам новый способ
доказательства элементарной эквивалентности.

\section{Понижение мощности}
        \label{lowenheim-skolem-1}%

В этом разделе мы опишем приём, позволяющей в интерпретации
большой мощности выделять часть, которая будет элементарно
эквивалентна исходной (и, более того, исходная будет её
элементарным расширением в смысле определения на
с.~\pageref{elementary-extension}). Например, во всяком
бесконечном упорядоченном множестве этот приём позволит найти
счётное подмножество, элементарно эквивалентное исходному как
интерпретация сигнатуры $({=},{<})$.

С помощью этой конструкции (составляющей содержание теоремы
Левенгейма\ч Сколема об элементарной подмодели) легко дать
обещанное на с.~\pageref{dense-sets-elementary-equivalence}
другое доказательство того, что любые два плотно упорядоченных
множества без первого и последнего элемента элементарно
эквивалентны. В самом деле, выберем в них счётные
части, элементарно эквивалентные целым множествам. Эти части
будут плотными и не будут иметь первого и последнего элементов,
так как эти свойства записываются формулами. Как известно (см.,
например,~\cite{naive-set-theory}), любые два счётных плотных
упорядоченных множества без первого и последнего элемента
изоморфны, и потому (теорема~\ref{isomorphism-equivalence},
с.~\pageref{isomorphism-equivalence}) элементарно эквивалентны.
Следовательно, и исходные множества будут элементарно
эквивалентны.

Прежде всего уточним слова \лк часть
интерпретации\пк. Если сигнатура состоит только из предикатных
символов, то проблем нет: взяв произвольное непустое
подмножество~$X$ произвольной интерпретации, мы можем ограничить
предикаты на~$X$ и получить новую интерпретацию. Если в
сигнатуре есть функциональные символы, мы должны ещё
потребовать, чтобы~$X$ было замкнуто относительно
соответствующих функций (значения функций на элементах
подмножества~$X$ должны лежать в~$X$). Возникающая при этом
интерпретация с носителем~$X$ называется \emph{подструктурой}
исходной.\index{Подструктура}

\begin{theorem}[Левенгейма\ч Сколема об элементарной подмодели]
\index{Теорема!Левенгейма\ч Сколема об элементарной
подмодели|defin}%
        \label{lowenheim-skolem-theorem}%
Пусть имеется конечная или счётная сигнатура~$\sigma$ и
некоторая бесконечная интерпретация $M$ этой сигнатуры. Тогда
можно указать счётное подмножество подмножество $M'\subset M$,
которое будет подструктурой $M$ (замкнуто относительно
сигнатурных функций) и для которого $M$ будет элементарным
расширением $M'$.
        %
\end{theorem}

\begin{proof}
        %
Начнём с первого требования теоремы: $M'$ должно быть подструктурой.
Само по себе его выполнить легко, как говорит следующая лемма.

\textsf{Лемма 1}. Пусть $A$\т произвольное конечное или
счётное подмножество множества~$M$. Тогда существует конечное или счётное
множество $B\subset M$, содержащее $A$, которое является подструктурой
(замкнуто относительно сигнатурных функций в $M$).

Утверждение леммы почти очевидно: надо добавить к $A$ результаты
применения всех функций к его элементам, потом результаты
применения всех функций к добавленным элементам и так счётное
число раз. (Другими словами, надо добавить значения всех термов
сигнатуры на оценках, при которых индивидные переменные
принимают значения в $A$.) Ясно, что получится конечное или
счётное множество, так как на каждом шаге расширения добавляется
счётное множество новых элементов и шагов счётное число. (Можно
заметить также, что термов счётное число.) Лемма~1 доказана.

Замкнутость подмножества $A$ множества $M$ относительно
сигнатурных функций позволяет рассматривать интерпретацию с
носителем~$A$ и с индуцированными из~$M$ функциями и
предикатами. Однако она, конечно, не обязана быть элементарно
эквивалентной $M$, как показывает множество очевидных примеров.
(Если, скажем, в сигнатуре нет функций, а одни предикаты,
то любое подмножество будет замкнуто.)

Поэтому нам необходимо ещё одно свойство замкнутости. Пусть
$A$\т некоторое подмножество $M$ (напомним, что мы рассматриваем
интерпретацию сигнатуры~$\sigma$ с носителем~$M$). Назовём~$A$
\emph{экзистенциально замкнутым},%
\index{Экзистенциально замкнутое множество|defin}%
\index{Множество!экзистенциально замкнутое|defin}
если для всякой
формулы $\varphi(x,x_1,\dots,x_n)$ сигнатуры~$\sigma$ и для
любых элементов $a_1,\dots,a_n\in A$ выполнено такое
утверждение: если существует $m\in M$, для которого (в $M$)
истинно $\varphi(m,a_1,\dots,a_n)$, то элемент $m$ с~таким
свойством можно выбрать и внутри~$A$.

(Более формально следовало бы сказать, что для всякой формулы
$\varphi$, параметры которой содержатся среди $x,x_1,\dots,x_n$,
и для любых элементов $a_1,\dots,a_n\hm\in A$ выполнено такое
утверждение: если существует $m\hm\in M$, для которого $\varphi$
истинна в интерпретации~$M$ на оценке $(x\hm\mapsto m,x_1\hm\mapsto
a_1,\dots,x_n\hm\mapsto a_n)$, то существует и
элемент $m\hm\in A$ с тем же свойством.)

Обратите внимание, что в этом определении (в отличие от
формулировки теоремы) не идёт речь об истинности какой бы то ни
было формулы в $A$\т только об истинности в $M$. В нём говорится
примерно вот что: если вообще (во всём $M$) найдётся элемент,
связанный неким формульным отношением с элементами
$a_1,\dots,a_n\in A$, то такой элемент найдётся и внутри
самого~$A$.

\textsf{Лемма 2}. Пусть $A\subset M$\т конечное или
счётное множество. Тогда существует конечное или счётное
множество $B\subset M$, содержащее $A$ и являющееся экзистенциально
замкнутым.

Доказательство леммы 2 аналогично доказательству предыдущей
леммы: формул $\varphi$ счётное число и конечных наборов
элементов из $A$ тоже счётное число. Поэтому можно посмотреть, в
каких случаях элемент $m$ из определения экзистенциальной
замкнутости существует, и добавить один из таких элементов
(здесь используется аксиома выбора). Один раз так сделать
недостаточно, так как добавленные элементы также могут
использоваться в качестве $a_1,\dots,a_n$ в определении, поэтому
такую процедуру надо повторить счётное число раз и взять
объединение полученных множеств. Оно уже будет экзистенциально
замкнуто (любой набор получается на конечном шаге и на следующем
шаге он обслуживается, если нужно). Лемма~2 доказана.

На самом деле леммы~1 и~2 можно соединить.

\textsf{Лемма 3}. Пусть $A\subset M$\т конечное или
счётное множество. Тогда существует конечное или счётное
множество $B\subset M$, содержащее $A$, замкнутое относительно
сигнатурных функций и экзистенциально замкнутое.

В самом деле, чтобы получить такое множество~$B$, достаточно
чередовать шаги замыкания относительно сигнатурных функций
и экзистенциального замыкания, а потом взять объединение
полученной последовательности множеств. Лемма~3 доказана.

\textsf{Лемма~4}. Пусть $M'\subset M$ замкнуто относительно
сигнатурных функций и экзистенциально замкнуто. Тогда $M$
является элементарным расширением $M'$.

Отсюда уже вытекает утверждение
теоремы~\ref{lowenheim-skolem-theorem}: применим лемму~3 к
некоторому счётному подмножеству множества $M$, а затем
воспользуемся леммой~4.

Доказательство леммы~4 также довольно просто. Напомним
определение элементарного расширения: требуется, чтобы
        $$
M'\vDash \varphi(a_1,\dots,a_n) \quad
\Leftrightarrow \quad
M \vDash \varphi(a_1,\dots,a_n)
        $$
для любой формулы $\varphi(x_1,\dots,x_n)$ и для любых~$a_1,\dots,a_n\hm\in M'$.

(Формально следовало бы сказать: для любой формулы с параметрами
и любой оценки, при которой все параметры принимают значения в
$M'$, истинность этой формулы в $M'$ на этой оценке равносильна
истинности той же формулы в $M$ на той же оценке.)

Будем доказывать это индукцией по построению формулы $\varphi$.
Для атомарных формул это очевидно: значения термов не
зависят от того, проводим ли мы вычисления в $M$ или $M'$, а
предикаты на $M'$ индуцированы из $M$.

Если формула $\varphi$ есть конъюнкция, дизъюнкция, импликация
или отрицание, то её истинность как в $M$, так и в $M'$ определяется
истинностью её частей (и можно сослаться на предположение индукции).

Единственный нетривиальный случай\т если формула $\varphi$
начинается с квантора. Мы можем сократить себе работу и
рассматривать только квантор существования, так как $\forall\xi$
можно заменить на $\lnot\exists\xi\lnot$. 

Итак, пусть формула
$\varphi(x_1,\dots,x_n)$ имеет вид
        $
\exists x\, \psi (x, x_1,\dots,x_n).
        $
Если $M'\vDash \varphi(a_1,\dots,a_n)$ для некоторых
$a_1,\dots,a_n\hm\in M'$, то найдётся
элемент $m\hm\in M'$, для которого
$M'\vDash\psi(m,a_1,\dots,a_n)$. Тогда по предположению индукции
(формула~$\psi$ короче формулы~$\varphi$) можно перейти к
большей интерпретации и заключить, что
$M\vDash\psi(m,a_1,\dots,a_n)$, и потому $M\vDash \varphi(a_1,\dots,a_n)$. Обратное
рассуждение просто так не проходит, поскольку существующий
элемент~$m$ существует в $M$, а не в $M'$, и предположение
индукции применить нельзя. Однако ровно для этого у нас есть
требование экзистенциальной замкнутости, которое позволяет
заменить элемент~$m$ на другой элемент из~$M'$ и завершить
доказательство.
        %
\end{proof}

Вот пример применения теоремы Левенгейма\ч Сколема в алгебре:
существует алгебраически замкнутое счётное подполе поля $\bbC$
комплексных чисел. (В самом деле, требование алгебраической замкнутости
можно записать в виде счётной последовательности формул\т по одной
для каждой степени многочлена. Аксиомы поля также можно записать
в виде формул. Значит, счётная элементарная подмодель
поля $\bbC$ будет также алгебраически замкнутым полем.)

Впрочем, алгебраистов такое применение скорее насмешит\т они и
так знают, что алгебраические элементы поля $\bbC$ (корни
многочленов с целыми коэффициентами) образуют счётное
алгебраически замкнутое поле.

Любопытный парадокс связан с попытками применить теорему
Левенгейма\ч Сколема в теории множеств. Представим себе
интерпретацию языка теории множеств (предикаты $=$ и $\in$),
носителем которой является множество всех множеств. Такого
множества, строго говоря, не бывает, но если про это забыть и
применить теорему Левенгейма\ч Сколема об элементарной
подмодели, то можно оставить лишь счётное число множеств так,
чтобы истинность утверждений теории множеств не изменилась. Но
среди этих утверждений есть и утверждение о существовании
несчётного множества\т как же так? Это рассуждение содержит
столько пробелов, что указать один из них совсем нетрудно. Тем
не менее оно может быть переведено в аксиоматическую теорию
множеств и даёт интересные (хотя уже не парадоксальные)
результаты.

Два дополнительных замечания усиливают теорему Левенгейма\ч
Сколема. Во\д первых, легко видеть, что для всякого конечного
или счётного подмножества $A\hm\subset M$ найдётся счётная
элементарная подструктура $M'\hm\subset M$, содержащая все
элементы $A$. (В самом деле, процесс замыкания, использованный
при доказательстве, можно начинать с множества $A$.)

Во\д вторых, можно отказаться от требования счётности сигнатуры
и сказать так: для всякого подмножества $A\subset M$ найдётся
элементарная подструктура $M'\subset M$, содержащая~$A$,
мощность которой не превосходит максимума из $\aleph_0$,
мощности множества~$A$ и мощности сигнатуры. В самом деле, и
конструкция замыкания относительно сигнатурных операций, и
конструкция экзистенциального замыкания, и счётное объединение
возрастающей цепи не выводят мощность за пределы указанного
максимума, поскольку и формулы, и термы являются конечными
последовательностями символов сигнатуры и счётного числа других
символов (см.~подробнее в~\cite{naive-set-theory}); то же самое
можно сказать о числе возможных наборов значений параметров.

Мы научились уменьшать мощность структуры, не меняя множества
истинных в ней формул. Можно, напротив, увеличивать мощность
(соответствующее утверждение иногда называют теоремой
Левенгейма\ч Сколема об элементарном расширении). Но эта
конструкция использует теорему компактности для языков первого
порядка, которая в свою очередь вытекает из теоремы Гёделя о
полноте. Поэтому мы отложим обсуждение этого утверждения до
следующей главы.
